\chapter{Dataset Preparation and Comparative Model Training}
\label{cap:segmentation}
\label{cap:dataset_training}

Following the theoretical analysis of deep learning architectures in Chapter~\ref{cap:architectures}, this chapter describes the datasets, annotation pipelines, and experimental model evaluations that guide the configuration for robotic integration.

The experimental study began by fine-tuning FCN, DeepLabV3, and LR-ASPP from pretrained weights on a manually annotated lettuce dataset for an empirical comparison. Because manual segmentation was time-consuming, this initial training was limited to lettuce rather than repeated across every target crop. Its outcome led to the adoption of LR-ASPP for embedded deployment.

The second experiment, carried out after the selection of LR-ASPP, tests whether networks trained on masks generated automatically by SAM~3, without any manual correction, can replace manual annotation. Manual masks were available for lettuce and were more accurate, but the aim was to evaluate labels that require no human work.

The experiment uses only lettuce images, because lettuce is the only crop with manual ground truth. SAM~3 generated three label sets for these images, with four, three, and two classes (Section~\ref{subsec:c5_semantic_classes}), and a separate LR-ASPP network was trained on each. The predictions of each network are compared with the manual lettuce masks on all its classes. The path class matters most, because it is the only one used for row following. The SAM~3 masks themselves are not compared with the manual masks.

The binary labeling was then extended to the complete image collection, labeled entirely by SAM~3 (\texttt{Data\_allPlantsSAM2}), and a network trained on this collection is the one deployed on the robot. The results of this configuration are reported in Table~\ref{tab:c5_results_binary_general}.

This chapter details image acquisition, dataset partitioning, training parameters, and offline benchmark evaluations. Subsequent chapters detail geometric navigation extraction and closed-loop robot validation.

\section{Image acquisition and class definition}
\label{sec:c5_acquisition_classes}

\subsection{Image sources and acquisition conditions}
\label{subsec:c5_image_acquisition}

\placeholderfigure{fig:crop_dataset_samples}%
{Sample field images across crops and corresponding training masks}%
{Esempi visivi delle acquisizioni di campo per lattuga, pomodoro e cetriolo, con le corrispondenti maschere di annotazione impiegate per l'addestramento (ground truth manuale per la lattuga, maschere binarie del path autogenerate da SAM~3 per pomodoro e cetriolo).}%
{5.5cm}

\todobox{Origine delle immagini e condizioni di acquisizione}{%
Specificare la sorgente dei dati e le modalit\`a di raccolta:
\begin{itemize}
  \item Dispositivo di ripresa utilizzato (risoluzione e ottica), quota di acquisizione e angolazione.
  \item Condizioni operative nei campi (aziende agricole o serre sperimentali, condizioni meteo e fasce orarie).
  \item Numero complessivo di immagini raccolte per lattuga, pomodoro e cetriolo.
\end{itemize}
}

\subsection{Semantic classes and labeling targets}
\label{subsec:c5_semantic_classes}

The segmentation labels in this thesis use up to four semantic classes: background, plants, path (the traversable corridor between two crop rows), and covering material (plastic mulch film). The manual lettuce annotations use all four classes (Section~\ref{subsec:c5_lettuce}).

The labels generated with SAM~3 (Section~\ref{subsec:c5_sam_annotation}) follow three configurations, each used to train a separate network:
\begin{itemize}
\item \textbf{Four-class configuration:} background, plants, path, and covering material, matching the four classes of the manual lettuce annotations.
\item \textbf{Three-class configuration:} background, plants, and path. No prompt targets the covering material.
\item \textbf{Binary configuration:} background and path. This configuration is used both for the lettuce-only binary network and for the multi-crop dataset \texttt{Data\_allPlantsSAM2}.
\end{itemize}
For row following, the path mask is the input of the geometric extraction in Chapter~\ref{cap:geometric_reference}.

\section{Dataset preparation}
\label{sec:c5_datasets}

\subsection{Lettuce dataset with manual annotation}
\label{subsec:c5_lettuce}

The manually annotated lettuce dataset, identified as \texttt{Data},
contains 762 image--mask pairs distributed across eight folders. Its
masks were prepared from the manual annotations stored in
\texttt{annotations.json}. The same manual reference is used for the
initial architecture comparison and for evaluating the models trained
with SAM~3-generated labels. These manual masks are the ground truth for
all evaluations in this chapter. They contain four classes (background,
plants, path, and covering material), and the path class marks only the
corridor the robot should follow, not other visible paths. The three lettuce autolabel datasets
(\texttt{Data\_insalataSAM\_4class}, \texttt{Data\_insalataSAM\_3class},
and \texttt{Data\_insalataSAM\_2class}) contain images from the same eight
lettuce folders, with masks generated by SAM~3 instead of the manual
annotations.

\subsection{Automated zero-shot annotation via SAM~3}
\label{subsec:c5_sam_annotation}

Manual segmentation requires substantial human effort, which motivated testing whether SAM~3 can replace it. The experiment was carried out without any manual work: SAM~3 was run offline with text prompts, and its masks were used for training without inspection, filtering, or correction. Manual masks were available for lettuce and were more accurate, but keeping the unedited masks was an intentional choice, to measure how well a network performs when the annotation step is fully automatic.

All autolabel masks were generated using \texttt{SAM3SemanticPredictor} with the local \texttt{sam3.pt} checkpoint and a confidence threshold of 0.25. For each text prompt, SAM~3 returned zero or more candidate masks, and pixels with mask values greater than 0.5 were assigned to the corresponding class. When masks of different prompts overlapped, the class was assigned according to a fixed prompt order, listed below for each dataset: later prompts were written only on pixels still labeled as background, so earlier classes were never overwritten. Pixels not assigned by any prompt remained background. The prompts and their order are:
\begin{itemize}
    \item \texttt{Data\_insalataSAM\_4class} (lettuce images): \texttt{plants} (class 1), then \texttt{plastic mulch film} (class 3), then \texttt{central farm path} (class 2)
    \item \texttt{Data\_insalataSAM\_3class} (lettuce images): \texttt{plants} (class 1), then \texttt{central farm path} (class 2)
    \item \texttt{Data\_insalataSAM\_2class} (lettuce images): \texttt{central farm path} only (class 1)
    \item \texttt{Data\_allPlantsSAM2} (lettuce, tomato, and cucumber images): \texttt{central path} only (class 1).
\end{itemize}
The SAM~3 masks are not compared directly with the manual annotations. Their quality enters the evaluation only through the networks trained on them (Section~\ref{subsec:c5_supervision_configurations}).

\placeholderfigure{fig:sam3_annotation_pipeline}%
{SAM~3 zero-shot autolabeling and mask composition workflow}%
{Diagramma della pipeline di autolabeling zero-shot: formulazione dei prompt testuali (\texttt{"plants"}, \texttt{"central farm path"}, \texttt{"plastic mulch film"}, \texttt{"central path"}), generazione automatica delle maschere tramite SAM~3 e composizione deterministica delle etichette senza filtraggio o revisione manuale.}%
{5cm}

\subsection{Tomato and cucumber images}
\label{subsec:c5_tomato_cucumber}

In addition to lettuce, the image collection includes tomato and cucumber crops. These images were annotated exclusively with SAM~3 and are included only in the multi-crop dataset (\texttt{Data\_allPlantsSAM2}) used to train the binary network deployed on the robot. They were not included in the training sets of the lettuce-specific models.

Because no manual ground truth was created for tomato and cucumber, these crops are not part of the quantitative evaluations in Section~\ref{sec:c5_supervision_comparison}. Similarly, the initial comparison among architectures (FCN, DeepLabV3, and LR-ASPP) was carried out exclusively on the lettuce dataset.

\subsection{Dataset partitioning}
\label{subsec:c5_dataset_splits}

The lettuce training runs use an 80\% training, 10\% validation, and
10\% test partition. For the 762 image--mask pairs, this yields 609 training, 76 validation, and 77
test samples.

The training split is used to update the network weights. The
validation split is used to compute the loss at the end of each epoch and
to select the best checkpoint. The test split is not used during training
and is used only afterwards, to compute the final scores. The three lettuce
autolabel datasets contain the same 762 images as \texttt{Data}, in the
same order, and the three autolabel projects use a split file identical
to that of the manual model. The 77 test images are therefore the same
physical images for all four lettuce networks, and none of them belongs
to the 609 training images of any of these networks. The multi-crop binary network evaluated in Section~\ref{sec:c5_supervision_comparison}
is trained on the complete image collection (\texttt{Data\_allPlantsSAM2},
with lettuce, tomato, and cucumber images). Before splitting this
collection, the 77 lettuce test images were removed from it, so that none
of them is used for training. The remaining 1882 of 1959 image--mask pairs
were then split 80/10/10.

The split is random at the level of individual images: no test image is
used for training, but each lettuce folder, which corresponds to one
acquisition video recorded along a different path, contributes images to
both the training and the test split. During training, the networks
therefore see other photos taken along the same paths as the test images. The reported IoU values may
overestimate performance on paths or fields not seen during training. This
limitation applies to all networks of the comparison.

The realistic scenario uses only the test images from \texttt{insalata\_grande1}, \texttt{insalata\_grande2}, and \texttt{insalata\_media1}: 42 of the 77 test images. These folders were recorded between parallel rows, in the area where the robot operates, from viewpoints the robot can reach. The other five folders were recorded from viewpoints the robot cannot reach and can show more than one path, while the manual masks label only the path the robot should follow. The realistic scenario is therefore a subset of the test split, not a separate test set, and its results describe only the robot's operating conditions.

\section{Architecture comparison: FCN, DeepLabV3, and LR-ASPP}
\label{sec:c5_architecture_comparison}

This section selects the segmentation architecture used in the rest of
the thesis. FCN, DeepLabV3, and LR-ASPP were trained and compared on the
manually annotated lettuce dataset.

\subsection{Model selection and pretrained weights}
\label{subsec:c5_model_selection}

FCN, DeepLabV3, and LR-ASPP were chosen from the implementations available
in TorchVision. Their selection was practical, rather than the outcome of a
comparison with every available architecture.

The implementation separates each network into a backbone, which extracts
image features, and a segmentation head, which produces class scores at
each output location. Transfer learning is implemented by loading
pretrained parameters and adapting the output layers to the target
labeling task, rather than initializing every network parameter from
scratch.

For FCN and DeepLabV3, the selected implementations use a ResNet-50
backbone with pretrained segmentation weights. After these weights
are loaded, the final convolution of the main classifier is replaced by
a new $1\times1$ convolution with one output channel per target class
(4 channels). The final convolution of the auxiliary classifier is
replaced in the same way.

LR-ASPP is initialized in the same way. The implementation loads the
default pretrained weights of the complete LR-ASPP model with a
MobileNetV3-Large backbone. Its segmentation head ends with two
$1\times1$ convolutions, one applied to the low-level features and one
to the high-level features~\cite{torchvision_lraspp}. Both are replaced
by new convolutions with 4 output channels. LR-ASPP has no auxiliary
classifier.

The default weights of the three models were trained on a subset of
COCO restricted to the 20 categories of Pascal VOC~\cite{torchvision_fcn,torchvision_deeplabv3,torchvision_lraspp}.
All three networks therefore start from pretrained segmentation weights
of the same origin, and only their final output convolutions are newly
initialized.

\subsection{Training configuration and loss formulation}
\label{subsec:c5_training_config}

The optimization pipeline addresses class imbalance and follows a two-stage transfer learning schedule for all evaluated architectures.

\subsubsection{Loss function and class weighting}
Agricultural scenes exhibit a strong pixel imbalance between the navigable soil corridor, the plant canopy, and background areas. To prevent the optimization from biasing toward dominant categories, the models are trained with a weighted cross-entropy loss, denoted $\mathcal{L}_{\text{WCE}}$, in which the contribution of each pixel is multiplied by the weight $w_c$ of its ground-truth class $c$, defined in Eq.~\eqref{eq:class_weights_norm}.

Class weights are computed exclusively on the training split $\mathcal{D}_{\text{train}}$ to avoid information leakage from validation and test partitions. For each class $c$, the total pixel count $N_c = \sum_{k \in \mathcal{D}_{\text{train}}} \sum_{i,j} \mathbb{I}(Y^{(k)}_{i,j} = c)$ is accumulated across all training masks. The balanced inverse frequency $\text{inv}_c$ is then normalized by the mean of all inverse frequencies:
\begin{equation}
    \text{inv}_c = \frac{\sum_{k=0}^{C-1} N_k}{C \cdot N_c}, \qquad w_c = \frac{\text{inv}_c}{\frac{1}{C}\sum_{k=0}^{C-1} \text{inv}_k}.
    \label{eq:class_weights_norm}
\end{equation}
This normalization enforces $\frac{1}{C} \sum_{c=0}^{C-1} w_c = 1.0$.

For FCN and DeepLabV3, which incorporate an auxiliary classifier on the \texttt{layer3} features (output stride 8 in the dilated TorchVision backbone), the total training objective combines the primary and auxiliary losses with an auxiliary weighting factor of $0.4$:
\begin{equation}
    \mathcal{L}_{\text{total}} = \mathcal{L}_{\text{WCE}}(\text{main}) + 0.4 \cdot \mathcal{L}_{\text{WCE}}(\text{aux}).
    \label{eq:total_training_loss}
\end{equation}
For LR-ASPP, which lacks an auxiliary head, training relies solely on the primary loss $\mathcal{L}_{\text{WCE}}(\text{main})$.

\subsubsection{Data augmentation and two-stage optimization}
Data augmentation is applied only to the training subset (\texttt{full\_\allowbreak train\_\allowbreak dataset}): random horizontal flips with probability 0.5, random rotations between $-10^\circ$ and $10^\circ$, and, with probability 0.8, a color jitter of brightness and contrast (factors in $[0.8, 1.2]$) and saturation (factors in $[0.85, 1.15]$) to handle variable outdoor illumination. The validation and test sets (\texttt{full\_\allowbreak eval\_\allowbreak dataset}) are not augmented. Images are not resized, and all images are normalized with the ImageNet mean and standard deviation.

The networks were trained on an NVIDIA GeForce RTX 2080 SUPER with 8~GB of memory, with a batch size of 2.

Optimization proceeds across $200$ epochs in two distinct phases:
\begin{enumerate}
    \item \textbf{Stage 1 (Warmup, Epochs 1--10):} Backbone parameters are frozen by setting \texttt{requires\_grad = False}. Only the segmentation heads (and auxiliary classifiers when present) are optimized using AdamW with an initial learning rate of $10^{-3}$ and weight decay of $10^{-4}$. A cosine-annealing scheduler with $T_{\text{max}} = 200$ and $\eta_{\text{min}} = 10^{-6}$ updates the learning rate after each epoch. In the first epochs, the newly initialized output layers make large errors, and the resulting large gradients could damage the pretrained weights. Freezing the reused layers during these epochs gives the new layers time to learn reasonable weights~\cite[Chapter~11]{geron2026handson}.
    \item \textbf{Stage 2 (Differential Fine-Tuning, Epochs 11--200):} At epoch 10, the backbone is unfrozen. A new AdamW optimizer is instantiated with decoupled parameter groups: the backbone is assigned a learning rate of $10^{-5}$ and the segmentation heads a learning rate of $10^{-4}$. The lower learning rate of the unfrozen backbone limits changes to its pretrained weights~\cite[Chapter~11]{geron2026handson}. The cosine-annealing scheduler is re-initialized across the remaining 190 epochs ($T_{\text{max}} = 190, \eta_{\text{min}} = 10^{-6}$).
\end{enumerate}

Early stopping was not used: all three networks were trained for the full 200 epochs, so that the complete training and validation loss curves could be recorded. After each epoch, the model is evaluated on the unaugmented validation set. The checkpoint \texttt{best.pth} is saved whenever the validation loss falls below its previous minimum, and \texttt{last.pth} stores the state at the final epoch.

\placeholderfigure{fig:training_loss_curves}%
{Training and validation loss across epochs for FCN, DeepLabV3, and LR-ASPP}%
{Curve di apprendimento durante le 200 epoche: loss di training e loss di validazione (dai file \texttt{training\_log.csv}) per le tre architetture candidate sul dataset di lattuga.}%
{4.5cm}

\subsection{Evaluation metrics and results}
\label{subsec:c5_architecture_results}

The segmentation scores come from the three
\texttt{*\_\allowbreak test\_\allowbreak metrics.csv} files produced by
the test script on the 77 test images, one per network. Each contains the
aggregate metrics, the number of parameters, the checkpoint epoch, and the
IoU values for background, plants, path, and covering material. The
reported mIoU agrees with the arithmetic mean of these four class IoUs,
including background. The same script records the inference time on
the NVIDIA GeForce RTX 2080 SUPER used for training. Each test image is
processed one at a time at its original resolution, and the measured
time covers the forward pass and the argmax. These times are not
measurements of the complete robot processing or control rate.

Table~\ref{tab:c5_architecture_results} reports the exported metrics.
The checkpoint epoch is the epoch of the evaluated \texttt{best.pth},
the checkpoint with the lowest validation loss. It is not the training
duration, since every network was trained for 200 epochs.

\begin{table}[htbp]
    \centering
    \small
    \caption{Architecture comparison on the 77 test images of the manual
    lettuce dataset. Parameter counts refer to the trained
    four-class models and include the auxiliary classifier of FCN and
    DeepLabV3, which is used only during training.}
    \label{tab:c5_architecture_results}
    \begin{tabular}{@{}l r r r r r r@{}}
        \toprule
        Model & \shortstack{Parameters\\(M)} & Epoch & mIoU & Path IoU &
        \shortstack{Mean time\\(ms)} & FPS \\
        \midrule
        FCN & 35.3 & 38 & 0.9253 & 0.9654 & 59.10 & 16.92 \\
        DeepLabV3 & 42.0 & 43 & 0.9253 & 0.9686 & 86.90 & 11.51 \\
        LR-ASPP & 3.2 & 162 & 0.9145 & 0.9532 & 8.21 & 121.87 \\
        \bottomrule
    \end{tabular}
\end{table}

FCN and DeepLabV3 reach the same mIoU on the test images (0.9253, with a
difference below $10^{-4}$), compared with 0.9145 for LR-ASPP, a deficit of
1.08 percentage points. Its mean recorded time is 8.21~ms, compared
with 59.10~ms and 86.90~ms. These timings correspond to a 7.20-fold
reduction relative to FCN and a 10.59-fold reduction relative to
DeepLabV3.

Table~\ref{tab:c5_architecture_class_iou} provides the class-level
comparison. LR-ASPP has lower IoU than both alternatives in each class,
although its path IoU remains 0.9532. The data support a trade-off
between accuracy and latency, not a claim of equal segmentation accuracy.

\begin{table}[htbp]
    \centering
    \small
    \caption{Per-class IoU on the same 77 test images.}
    \label{tab:c5_architecture_class_iou}
    \begin{tabular}{@{}l r r r r@{}}
        \toprule
        Model & Background & Plants & Path & Covering material \\
        \midrule
        FCN & 0.9493 & 0.9385 & 0.9654 & 0.8479 \\
        DeepLabV3 & 0.9494 & 0.9384 & 0.9686 & 0.8448 \\
        LR-ASPP & 0.9439 & 0.9262 & 0.9532 & 0.8346 \\
        \bottomrule
    \end{tabular}
\end{table}

\subsection{Selection of LR-ASPP}
\label{subsec:c5_lraspp_selection}

LR-ASPP was selected mainly for its speed. Its mean inference time is
7.20 times lower than that of FCN and 10.59 times lower than that of
DeepLabV3, at the cost of an mIoU 1.08 percentage points lower. It was
not the most accurate model, but it offered the accuracy--latency
compromise sought for robot perception. The subsequent SAM-supervised
training runs therefore used LR-ASPP while varying the output label set.

The inference times have one limitation. The three networks were timed
in separate runs of the test script, and the load on the GPU during each
run was not recorded. These conditions affect the measured time: in the
evaluation scripts of Section~\ref{sec:c5_supervision_comparison}, an
LR-ASPP checkpoint records mean times between 0.010~s and 3.67~s per
image. The ratios of 7.20 and 10.59 are therefore reliable only if the
three runs had similar conditions.

These times were measured on the training PC, not on the robot. Whether
the complete perception and control pipeline on the Jetson Orin Nano
meets the control timing requirements has to be verified on the robot.

\section{Supervision strategy comparison}
\label{sec:c5_supervision_comparison}

\subsection{Training configurations and evaluation targets}
\label{subsec:c5_supervision_configurations}

Table~\ref{tab:c5_supervision_configurations} summarizes the five
LR-ASPP networks considered in this comparison. This study is separate
from the initial comparison of the three candidate architectures. The
first four networks are trained on lettuce images only, so that their
predictions can be evaluated against the manual ground truth of the same
crop. The fifth network reproduces the multi-crop binary configuration
deployed on the robot, trained on the complete SAM~3 collection without
the 77 lettuce test images.

\begin{table}[htbp]
    \centering
    \small
    \caption{LR-ASPP training configurations in the supervision comparison. Class counts include background; all networks are evaluated against the manual lettuce ground truth. The checkpoint epoch is the epoch of the evaluated \texttt{best.pth}, selected by minimum validation loss.}
    \label{tab:c5_supervision_configurations}
    \begin{tabularx}{\textwidth}{@{}l X X c r@{}}
        \toprule
        Configuration & Training data (labels) & Output classes & Classes & Epoch \\
        \midrule
        Manual baseline & \texttt{Data} (manual) & Background, plants, path, covering material & 4 & 162 \\
        SAM~3 four-class & \texttt{Data\_\allowbreak insalataSAM\_\allowbreak 4class} (SAM~3) & Background, plants, path, covering material & 4 & 30 \\
        SAM~3 three-class & \texttt{Data\_\allowbreak insalataSAM\_\allowbreak 3class} (SAM~3) & Background, plants, path & 3 & 38 \\
        SAM~3 binary & \texttt{Data\_\allowbreak insalataSAM\_\allowbreak 2class} (SAM~3) & Background, path & 2 & 18 \\
        SAM~3 binary, multi-crop & \texttt{Data\_\allowbreak allPlantsSAM2} without the 77 lettuce test images (SAM~3) & Background, path & 2 & 27 \\
        \bottomrule
    \end{tabularx}
\end{table}

Manual masks from \texttt{Data} supply the ground truth for all comparative evaluations. The evaluation measures the generalization capability of the trained networks rather than the static quality of the raw SAM~3 masks. Specifically, no direct pixel-level comparison is performed between uncurated SAM~3 masks and manual annotations; instead, each trained LR-ASPP network executes forward inference in evaluation mode (\texttt{eval()}), computing the discrete predicted class map via:
\begin{equation}
    \hat{Y}_{i,j} = \arg\max_{c \in \{0,\dots,C-1\}} P(Y_{i,j} = c \mid I),
    \label{eq:inference_argmax}
\end{equation}
which is subsequently compared against the human ground-truth mask.

Evaluation uses the Intersection over Union (IoU) metric:
\begin{equation}
    \text{IoU}_c = \frac{\text{TP}_c}{\text{TP}_c + \text{FP}_c + \text{FN}_c},
    \label{eq:eval_iou}
\end{equation}
where $\text{TP}_c$, $\text{FP}_c$, and $\text{FN}_c$ denote true positives, false positives, and false negatives for class $c$, respectively. Mean IoU (mIoU) denotes the arithmetic mean across all evaluated classes, including background where applicable.

Performance is evaluated on two subsets of the 77-image test partition:
\begin{itemize}
    \item \textbf{Complete test set (77 images):} Spans all eight lettuce acquisition folders (\texttt{insalata\_\allowbreak grande1}, \texttt{insalata\_\allowbreak grande2}, \texttt{insalata\_\allowbreak media1}, \texttt{insalata\_\allowbreak media2}, \texttt{insalata\_\allowbreak media3}, \texttt{insalata\_\allowbreak piccola}, \texttt{insalata\_\allowbreak piccola\_1}, and \texttt{insalata\_\allowbreak piccola\_2}), including viewpoints that the robot cannot reach.
    \item \textbf{Realistic operational scenario (42 images):} Restricts evaluation to \texttt{insalata\_\allowbreak grande1}, \texttt{insalata\_\allowbreak grande2}, and \texttt{insalata\_\allowbreak media1}, described in Section~\ref{subsec:c5_dataset_splits}. The networks are not retrained on this subset.
\end{itemize}

\subsection{Benchmark results across supervision strategies}
\label{subsec:c5_supervision_results}

All results in this section are computed on the test split. The four comparisons evaluated in \texttt{compareInsalata} are reported separately, each on the complete test set and on the realistic scenario.

\subsubsection{Three-class autolabel model: plants and path}
This comparison evaluates the two foreground categories shared by the 4-class manual baseline and the 3-class autolabel model (\texttt{InsalataSAM\_3class}). The other classes are ignored. As reported in Table~\ref{tab:c5_results_2class}, the manual model achieves a mean IoU of 0.9397 on the complete test set and 0.9536 in the realistic scenario. The 3-class autolabel model reaches a path IoU of only 0.4124 on the complete test set, and 0.8966 (plants IoU 0.8779) in the realistic scenario.

\begin{table}[htbp]
    \centering
    \small
    \caption{Test set performance of the three-class autolabel model on plants and path (records from \texttt{3\_classi} and \texttt{3\_classi\_realisticScenario}). Mean IoU denotes the arithmetic average of plants and path IoU.}
    \label{tab:c5_results_2class}
    \begin{tabularx}{\textwidth}{@{}l X r r r@{}}
        \toprule
        Scenario & Model & Plants IoU & Path IoU & Mean IoU \\
        \midrule
        \multirow{2}{*}{Complete test set} 
        & Manual baseline (4-class) & 0.9262 & 0.9532 & 0.9397 \\
        & Autolabel SAM (3-class) & 0.8130 & 0.4124 & 0.6127 \\
        \midrule
        \multirow{2}{*}{Realistic scenario} 
        & Manual baseline (4-class) & 0.9485 & 0.9587 & 0.9536 \\
        & Autolabel SAM (3-class) & 0.8779 & 0.8966 & 0.8873 \\
        \bottomrule
    \end{tabularx}
\end{table}

\subsubsection{Four-class autolabel model}
Table~\ref{tab:c5_results_4class} presents the complete 4-class evaluation comparing the manual baseline with the 4-class autolabel network (\texttt{InsalataSAM\_4class}). The manual model reaches an mIoU of 0.9145 on the complete test set and 0.9113 in the realistic scenario. The 4-class autolabel network reaches 0.4639 and 0.5925. In the realistic scenario the deficit is concentrated in the covering material (IoU 0.2495) and the background (IoU 0.3449), while its path IoU is 0.8997. Visual inspection of the SAM~3 masks showed that the \texttt{"plastic mulch film"} prompt segmented the covering poorly and also assigned the greenhouse plastic cover to this class. Since the greenhouse cover is background in the manual masks, this is consistent with the low IoU of both covering and background. A further plausible cause is that the prompt produces inaccurate masks where the film is soiled or damaged. These explanations come from visual inspection. The SAM~3 masks were not compared quantitatively with the manual annotations.

\begin{table}[htbp]
    \centering
    \small
    \caption{Complete four-class test set results (records from \texttt{4\_classi} and \texttt{4\_class\_realisticScenario}). mIoU is the average of all four classes including background.}
    \label{tab:c5_results_4class}
    \begin{tabularx}{\textwidth}{@{}l X r r r r r@{}}
        \toprule
        Scenario & Model & Background & Plants & Path & Covering material & mIoU \\
        \midrule
        \multirow{2}{*}{Complete test} 
        & Manual baseline & 0.9439 & 0.9262 & 0.9532 & 0.8346 & 0.9145 \\
        & Autolabel SAM & 0.4631 & 0.8081 & 0.4124 & 0.1719 & 0.4639 \\
        \midrule
        \multirow{2}{*}{Realistic test} 
        & Manual baseline & 0.9203 & 0.9485 & 0.9587 & 0.8178 & 0.9113 \\
        & Autolabel SAM & 0.3449 & 0.8757 & 0.8997 & 0.2495 & 0.5925 \\
        \bottomrule
    \end{tabularx}
\end{table}

\subsubsection{Binary autolabel model trained on lettuce only}
Table~\ref{tab:c5_results_binary_lettuce} compares path segmentation between the manual baseline and the binary autolabel network trained only on lettuce (\texttt{InsalataSAM\_2class}, records from \texttt{2\_classi} and \texttt{2\_classi\_realisticScenario}). Its labels come from the same \texttt{central farm path} prompt used for the path class of the 3-class and 4-class datasets. The network reaches a path IoU of 0.3917 on the complete test set and 0.9013 in the realistic scenario.

\begin{table}[htbp]
    \centering
    \small
    \caption{Path IoU of the binary autolabel model trained on lettuce only, on the test set. The manual model path is class~2; the binary model path is class~1.}
    \label{tab:c5_results_binary_lettuce}
    \begin{tabularx}{\textwidth}{@{}l X r@{}}
        \toprule
        Scenario & Model & Path IoU \\
        \midrule
        \multirow{2}{*}{Complete test set}
        & Manual baseline (4-class) & 0.9532 \\
        & Autolabel SAM (2-class, lettuce only) & 0.3917 \\
        \midrule
        \multirow{2}{*}{Realistic scenario}
        & Manual baseline (4-class) & 0.9587 \\
        & Autolabel SAM (2-class, lettuce only) & 0.9013 \\
        \bottomrule
    \end{tabularx}
\end{table}

\subsubsection{Multi-crop binary autolabel model}
Table~\ref{tab:c5_results_binary_general} reports the binary network trained on the multi-crop SAM~3 collection without the lettuce test images (project \texttt{allPlantsSAM\_\allowbreak test}, records from \texttt{2\_class\_\allowbreak allPlants\_\allowbreak clean} and its realistic-scenario variant). Its labels come from the \texttt{central path} prompt, and its training data contain tomato and cucumber images together with the remaining lettuce frames. The evaluated checkpoint is the one with the lowest validation loss (epoch 27, loss 0.1159). The network reaches a path IoU of 0.6369 on the complete test set and 0.9071 in the realistic scenario.

\begin{table}[htbp]
    \centering
    \small
    \caption{Path IoU of the multi-crop binary autolabel model, trained without the lettuce test images, on the test set. The manual model path is class~2; the binary model path is class~1.}
    \label{tab:c5_results_binary_general}
    \begin{tabularx}{\textwidth}{@{}l X r@{}}
        \toprule
        Scenario & Model & Path IoU \\
        \midrule
        \multirow{2}{*}{Complete test set}
        & Manual baseline (4-class) & 0.9532 \\
        & Autolabel SAM (2-class, multi-crop) & 0.6369 \\
        \midrule
        \multirow{2}{*}{Realistic scenario}
        & Manual baseline (4-class) & 0.9587 \\
        & Autolabel SAM (2-class, multi-crop) & 0.9071 \\
        \bottomrule
    \end{tabularx}
\end{table}

\subsection{Path segmentation with autolabel supervision}
\label{subsec:c5_binary_selection}

Table~\ref{tab:c5_path_summary} collects the path IoU of every network from Tables~\ref{tab:c5_results_2class}--\ref{tab:c5_results_binary_general}.

\begin{table}[htbp]
    \centering
    \small
    \caption{Path IoU on the test set for all networks of the supervision comparison. The multi-crop network is trained without the 77 lettuce test images.}
    \label{tab:c5_path_summary}
    \begin{tabular}{@{}l r r@{}}
        \toprule
        Network & Complete test set & Realistic scenario \\
        \midrule
        Manual baseline (4-class) & 0.9532 & 0.9587 \\
        SAM~3 four-class (lettuce) & 0.4124 & 0.8997 \\
        SAM~3 three-class (lettuce) & 0.4124 & 0.8966 \\
        SAM~3 binary (lettuce) & 0.3917 & 0.9013 \\
        SAM~3 binary (multi-crop) & 0.6369 & 0.9071 \\
        \bottomrule
    \end{tabular}
\end{table}

The results support the following conclusions:
\begin{enumerate}
    \item \textbf{The manual model is the best overall.} It has the highest IoU in every comparison, with path IoU 0.9532 on the complete test set and 0.9587 in the realistic scenario.
    \item \textbf{Lettuce-only autolabel training.} In the realistic scenario, the binary network trained only on lettuce reaches a path IoU of 0.9013, 5.7 percentage points below the manual model. The three lettuce autolabel networks differ by less than 0.5 percentage points on path IoU in this scenario (0.8966--0.9013), and on the complete test set the binary network is the lowest of the three (0.3917 against 0.4124). These data do not show that removing the plant and covering classes improves path segmentation.
    \item \textbf{Multi-crop autolabel data.} The multi-crop binary network reaches a path IoU of 0.9071 in the realistic scenario, close to the lettuce binary network (0.9013), and 0.6369 on the complete test set, against 0.3917. Its training data differ in several ways at once (tomato and cucumber images, a different set of lettuce frames, and the \texttt{central path} prompt), so the improvement cannot be attributed to one factor.
    \item \textbf{Complete test set versus realistic scenario.} Both scenarios evaluate the same networks, trained on all eight lettuce folders; the realistic scenario is only a subset of the test images. All autolabel networks lose much of their path accuracy on the complete test set (path IoU 0.39--0.41 for the lettuce networks and 0.64 for the multi-crop network), while the manual model does not (0.9532). The drop comes from the five folders outside the operational scenario. Their images were taken from viewpoints the robot cannot reach and can show several paths: the manual masks label only the path the robot should follow, whereas the autolabel networks segment every visible path, and the IoU counts the additional paths as false positives. In addition, visual inspection showed that SAM~3 produced less accurate path masks in images taken from unusual viewpoints or where the path between the rows is not clearly visible, so the networks were trained on noisier labels for these conditions. The complete-test values therefore combine segmentation errors, label noise, and this difference in label definition; the three effects were not separated quantitatively.
    \item \textbf{Full semantic segmentation.} With an mIoU of 0.5925 in the realistic scenario, the 4-class autolabel network does not support replacing manual annotation when all four classes are required.
    \item \textbf{Sufficiency for navigation.} Path IoU values of 0.9013 (lettuce binary) and 0.9071 (multi-crop binary) in the realistic scenario imply that more than 90\% of the ground-truth path pixels are predicted as path, since recall $\text{TP}/(\text{TP}+\text{FN})$ is never lower than IoU. Whether this accuracy is sufficient for row following depends on how the residual errors are distributed along the corridor boundaries. It must be established through an application threshold or a downstream test of the path extraction and controller (Chapters~\ref{cap:geometric_reference} and~\ref{cap:experiments}), not from IoU alone.
\end{enumerate}


\subsection{Network deployed on the robot}
\label{subsec:c5_deployed_network}

After the lettuce study, the binary labeling was extended to the complete image collection, labeled entirely by SAM~3 with the \texttt{central path} prompt and without manual correction (\texttt{Data\_allPlantsSAM2}). The binary LR-ASPP network trained on this collection (checkpoint \texttt{best\_allsam2.pth}, epoch 25) is the network used on the robot in Chapters~\ref{cap:platform} and~\ref{cap:experiments}. Its training data include the lettuce test images, so its own IoU on them would not be an independent measure and is not reported. The multi-crop network of Table~\ref{tab:c5_results_binary_general}, trained with the same labels on the collection without those images, gives the independent estimate for this configuration: path IoU 0.9071 in the realistic scenario and 0.6369 on the complete test set. It is a separate training run (best epoch 27 instead of 25), so the estimate refers to the configuration rather than to the exact deployed weights. The deployed network was also assessed qualitatively, by visual inspection of its predictions over the dataset, and then in simulation and in field tests (Chapter~\ref{cap:experiments}). The visual inspection also showed that it segments adjacent or non-traversable corridors, which the manual lettuce masks leave as background. The checkpoint documentation identifies \texttt{Data\_allPlantsSAM2} as its training set, whereas the historical training script refers to \texttt{Data\_allPlantsMIX}; the checkpoint does not store the dataset path.

For integration with the geometric extraction, a binary network outputs the path class directly. The extractor \texttt{pf4.py} receives the predicted class map and a \texttt{path\_class} parameter (default 1) and binarizes the map internally, so a multi-class output could also be used by setting \texttt{path\_class} to the index of its path class. It uses only the path class: it analyzes 24 horizontal strips and fits the left and right corridor edges with iteratively reweighted least squares (Chapter~\ref{cap:geometric_reference}).
